\section{Paper Description and Justification}
\label{sec:paper}

Relation extraction (RE) identifies the semantic relation between two entities in text. Supervised RE needs labelled data for every relation of interest, and most zero-shot methods still need labelled data for other, ``seen'' relations. \citet{zhou2024grasping} study a harder setting, \emph{definition-only zero-shot RE}: the only inputs are a natural-language definition of the target relation and an unlabelled corpus. Negative relations are assumed unknown.

Their framework, REPaL, has three stages. (1) \emph{Definition-based seed construction}: an LLM writes positive examples from the definition, and initial negatives are sampled from the unlabelled corpus. (2) \emph{Pattern learning}: a small language model (SLM), here the NLI model \texttt{roberta-large-mnli}, derived from RoBERTa \citep{liu2019roberta}, is fine-tuned as a binary classifier in which the definition is the hypothesis and the sentence is the premise. (3) \emph{Feedback}: the SLM's predictions on the unlabelled corpus, together with the synthesis history, are returned to the LLM, which writes follow-up positive examples and negative relation definitions targeted at the SLM's errors. The SLM is then retrained.

Evaluation resembles cross-validation. With $R$ test relations there are $R$ binary tasks, each treating one relation as positive and the rest as negative. Precision, recall and F1 are computed per task and averaged. Because F1 is averaged per relation, it can be lower than both mean precision and mean recall.

The paper reports REPaL (GPT-4o) at P/R/F1 of 78.86/77.28/74.61 on DefOn-FewRel and 68.98/47.63/47.80 on DefOn-Wiki-ZSL, ahead of the zero-shot RoBERTa NLI baseline (F1 55.07 and 35.28). Its ablation shows that the LLM's initial positives contribute most: removing them gives F1 55.07 on DefOn-FewRel, while removing positive feedback gives 70.03 and removing negative feedback gives 73.04.

\paragraph{Justification.} The work was published at EMNLP 2024 (pp.\ 13462--13486) by authors from the University of Illinois Urbana-Champaign and the University of Virginia. EMNLP is ranked A* in the ICORE2026 ranking,\footnote{\url{https://portal.core.edu.au/conf-ranks/}, checked 2026-11-09.} and the paper had 30 citations on Semantic Scholar when checked on 2026-11-09. Feasibility was strong: code is public, and the authors released the evaluation splits and a cache of GPT-4o synthesis, so replication needs no API spending. Scientifically, a method that needs only a definition should in principle transfer to new relation schemas, which is what this project's extension tests.
